\documentclass[../libro]{subfiles}

\begin{document}

\ifSubfilesClassLoaded{\mainmatter\chapter{行列式}\clearpage}{}

\KunAsteriskoEnEnhavtabelo
\section{\esperanto{La lasta leciono}}
\SenAsteriskoEnEnhavtabelo

\maldevigalegajxo

这是最后一课, \esperanto{la lasta leciono}.

这是好的!
您学完了本章.
您可能会想:
``我还能再学些什么?''
我认为, 这个问题, 难回答:
行列式只是一个工具.
不过, 我想,
您可以进一步地学习线性代数.

若您想见别的文献,
以下的说明或许对您是有用的.

习惯地, 几乎每一个 (说汉语的) 作者都说 ``矩阵'',
而不是 ``阵''
(当然, ``方阵'' 是一个例外).

习惯地, 人们叫一个 \(1 \times n\)~阵为一个行向量,
且叫一个 \(m \times 1\)~阵为一个列向量.
行向量与列向量可被统一地叫作向量.
% 不过, 当初我写作时, 不想给您多的挑战,
% 故我未正式地引入 ``向量'' 二字.

我一般表
\(A\)~的 \((i, j)\)-元%
以 \([A]_{i,j}\).
此记号自然是准确的;
并且,
我们可自由地代括号里的单文字 \(A\)
以复杂的文字, 如 \(3A + 4B\)
(当 \(A\), \(B\) 是同尺寸的阵时),
且方便地作计算:
\begin{align*}
    [3A + 4B]_{i,j}
        = [3A]_{i,j} + [4B]_{i,j}
        = 3[A]_{i,j} + 4[B]_{i,j}.
\end{align*}

其实, 习惯地,
多的作者不用记号 \([A]_{i,j}\),
而用较简单的记号 \(a_{ij}\)
(注意, 这儿没有逗号):
\begin{align*}
    \begin{bmatrix}
        a_{11} & a_{12} & \cdots & a_{1n} \\
        a_{21} & a_{22} & \cdots & a_{2n} \\
        \vdots & \vdots & {}     & \vdots \\
        a_{m1} & a_{m2} & \cdots & a_{mn} \\
    \end{bmatrix}
    =
    \begin{bmatrix}
        [A]_{1,1} & [A]_{1,2} & \cdots & [A]_{1,n} \\
        [A]_{2,1} & [A]_{2,2} & \cdots & [A]_{2,n} \\
        \vdots    & \vdots    & {}     & \vdots    \\
        [A]_{m,1} & [A]_{m,2} & \cdots & [A]_{m,n} \\
    \end{bmatrix}.
\end{align*}
于是, 在此写法下,
\(A\)~的 \((2, 3)\)-元是 \(a_{23}\),
且 \(B\)~的 \((4, 2)\)-元是 \(b_{42}\).
不过, 若 \(i\), \(j\) 的一个不能被单文字表示,
则逗号仍被用:
比如, \(C\)~的 \((11, 9)\)-元是 \(c_{11,9}\).
并且, 这个记号, 要求阵被单文字表示.
比如, 设 \(A\), \(B\) 分别是
\(m \times s\), \(s \times n\)~阵.
习惯地, 不写 \(AB\)~的 \((i, j)\)-元为
\({ab}_{ij}\)
(因为, \({ab}_{ij}\),
习惯地,
会被理解为%
一个数~\(a\) 跟
\(B\)~的 \((i, j)\)-元 \(b_{ij}\)
的积),
而记 \(C = AB\)
(也就是, 为 \(AB\) 取别名 \(C\)),
再说 \(AB\)~的 \((i, j)\)-元
\begin{align*}
    c_{ij}
    = a_{i1} b_{1j} + a_{i2} b_{2j} + \dots
    + a_{is} b_{sj}.
\end{align*}

一些作者不用括号 \({[} \ {]}\),
而用括号 \(( \ )\),
以包围数表, 作成一个阵:
\begin{align*}
    \begin{pmatrix}
        a_{11} & a_{12} & \cdots & a_{1n} \\
        a_{21} & a_{22} & \cdots & a_{2n} \\
        \vdots & \vdots & {}     & \vdots \\
        a_{m1} & a_{m2} & \cdots & a_{mn} \\
    \end{pmatrix}
    =
    \begin{bmatrix}
        a_{11} & a_{12} & \cdots & a_{1n} \\
        a_{21} & a_{22} & \cdots & a_{2n} \\
        \vdots & \vdots & {}     & \vdots \\
        a_{m1} & a_{m2} & \cdots & a_{mn} \\
    \end{bmatrix}.
\end{align*}

一些作者用大写拉丁字母 \(O\) 表示一个零阵.

一些作者用大写拉丁字母 \(E\) 表示一个单位阵.
一些作者用数字 \(1\) 表示一个单位阵.

一些作者用加粗的文字表示阵.
具体地, 他们写
\(\mathbf{A}\),
\(\mathbf{B}\),
\(\mathbf{C}\),
\(\mathbf{E}\),
\(\mathbf{I}\),
\(\mathbf{1}\),
\(\mathbf{O}\),
\(\mathbf{0}\),
\(\mathbf{x}\),
\(\dots\),
而不是
\(A\), \(B\), \(C\),
\(E\), \(I\), \(1\),
\(O\), \(0\),
\(x\), \(\dots\).

一些作者, 习惯地, 写%
一个方阵~\(A\) 的行列式
\(\det {(A)}\) 为 \(|A|\).
若 \(A\)~的元被具体地写出,
则有记号
\begin{align*}
    \begin{vmatrix}
        a_{11} & a_{12} & \cdots & a_{1n} \\
        a_{21} & a_{22} & \cdots & a_{2n} \\
        \vdots & \vdots & {}     & \vdots \\
        a_{n1} & a_{n2} & \cdots & a_{nn} \\
    \end{vmatrix}
    = {} &
    \left|\,
    \begin{bmatrix}
        a_{11} & a_{12} & \cdots & a_{1n} \\
        a_{21} & a_{22} & \cdots & a_{2n} \\
        \vdots & \vdots & {}     & \vdots \\
        a_{n1} & a_{n2} & \cdots & a_{nn} \\
    \end{bmatrix}
    \,\right|
    \\
    = {} &
    \det {
        \begin{bmatrix}
            a_{11} & a_{12} & \cdots & a_{1n} \\
            a_{21} & a_{22} & \cdots & a_{2n} \\
            \vdots & \vdots & {}     & \vdots \\
            a_{n1} & a_{n2} & \cdots & a_{nn} \\
        \end{bmatrix}
    }.
\end{align*}
其实, 行列式比阵先出现;
于是, 行列式有专门的竖线记号.

一些作者讲行列式时, 会提余子式与%
代数余子式.
通俗地, 一个方阵 \(A\)~的 \((i, j)\)-元
\(a_{i,j}\) 的%
余子式 \(M_{i,j}\),
% (这里, 为尽可能地消除歧义, 我用了逗号),
是行列式
\begin{align*}
    \begin{vmatrix}
        a_{1,1}     & \cdots & a_{1,j-1}   &
        a_{1,j+1}   & \cdots & a_{1,n}       \\
        \vdots      & {}     & \vdots      &
        \vdots      & {}     & \vdots        \\
        a_{i-1,1}   & \cdots & a_{i-1,j-1} &
        a_{i-1,j+1} & \cdots & a_{i-1,n}     \\
        a_{i+1,1}   & \cdots & a_{i+1,j-1} &
        a_{i+1,j+1} & \cdots & a_{i+1,n}     \\
        \vdots      & {}     & \vdots      &
        \vdots      & {}     & \vdots        \\
        a_{n,1}     & \cdots & a_{n,j-1}   &
        a_{n,j+1}   & \cdots & a_{n,n}       \\
    \end{vmatrix}.
\end{align*}
一个方阵 \(A\)~的 \((i, j)\)-元
\(a_{i,j}\) 的%
代数余子式 \(A_{i,j}\)
(其也被一些作者记为 \(C_{i,j}\)),
是
\((-1)^{i+j} M_{i,j}\).
于是,
对任何不超过 \(n\)~的正整数 \(i\), \(j\),
\begin{align*}
    |A|
    = {} &
    (-1)^{i+1} a_{i,1} M_{i,1}
    + (-1)^{i+2} a_{i,2} M_{i,2}
    + \dots
    + (-1)^{i+n} a_{i,n} M_{i,n}
    \\
    = {} &
    (-1)^{1+j} a_{1,j} M_{1,j}
    + (-1)^{2+j} a_{2,j} M_{2,j}
    + \dots
    + (-1)^{n+j} a_{n,j} M_{n,j},
\end{align*}
或
\begin{align*}
    |A|
    = {} &
    a_{i,1} A_{i,1}
    + a_{i,2} A_{i,2}
    + \dots
    + a_{i,n} A_{i,n}
    \\
    = {} &
    a_{1,j} A_{1,j}
    + a_{2,j} A_{2,j}
    + \dots
    + a_{n,j} A_{n,j}.
\end{align*}
我没有讲这二个概念,
因为我引入了子阵与它的记号
(一些作者不介绍表示子阵的记号).
其实,
余子式就是子阵的行列式:
\(
M_{i,j} = \det {(A(i|j))}
\).
这么看来, 我不必引入余子式或代数余子式了.

一些作者写方阵 \(A\)~的古伴
\(\operatorname{adj} {(A)}\)
为 \(A^{\ast}\),
且叫 ``古伴'' (即 ``古典伴随阵'')
为 ``伴随'' ``伴随阵'' ``伴随矩阵''.

我希望, 这些说明,
可助您适应其他的文献上的词与记号.
% 或许, 就像嬴政统一文字那样,
% 数学的符号也应被统一;
% 可是, 此事是难的.
% 我能作的事, 至多是助您适应这些符号.

% 最后, 我想说,
% 行列式并不只是代数的工具.
% 行列式在微积分与几何里也是有用的;
% 您可以去相关的文献里找到更多的讨论.
% (也可以直接去互联网看).
% 不过, 我对此了解不深;
% 再者, 若我进一步讲下去, 我就要
% ``一台电脑, 一个键盘, 一个图画一天''
% 了
% (乳胶可不是什么简单的东西,
% 但因为, 跟词儿比,
% 乳胶对我, 用鼠标有些困难的人,
% 更友好,
% 我忍了).

我就说这么多.
再见.

\end{document}

This is the last lesson;
this is the last section of chapter 1;
this is optional;
this is provided for those
who wish to know more about determinants
from other sources.
